% !TEX root = main.tex

%%%%%%%%%%%%%%%%%%%%%%%%%%%%%%%%%%%%%%%%%%%%%%%%%%%%%%
\section{図表}
%%%%%%%%%%%%%%%%%%%%%%%%%%%%%%%%%%%%%%%%%%%%%%%%%%%%%%
\subsection{図}
%%%%%%%%%%%%%%%%%%%%%%%%%%%%%%%%%%%%%%%%%%%%%%%%%%%%%%
　\texttt{pdf} ファイルを取り込むことができます．
% /////////////////////////////////////////////////////
\begin{figure}[H]
\centering
	\includegraphics[scale=0.5]{logo.pdf}
	\caption{舞鶴高専ロゴ}
	\label{fig:fig01}
\end{figure}
% /////////////////////////////////////////////////////
\textbf{図~\ref{fig:fig01}} のように，図番号を参照することもできます．

%%%%%%%%%%%%%%%%%%%%%%%%%%%%%%%%%%%%%%%%%%%%%%%%%%%%%%
\subsection{表}
%%%%%%%%%%%%%%%%%%%%%%%%%%%%%%%%%%%%%%%%%%%%%%%%%%%%%%
　表の作成は以下の通りです．
% /////////////////////////////////////////////////////
\begin{table}[H]
\centering
\caption{PID パラメータ}
\label{table:table01}
\begin{tabular}{|c|c|c|c|c|}
	\hline
	　 &
	比例感度 \({k}_{P}\) &
	比例帯 \({P}_{B}\) &	
	積分時間 \({T}_{I}\) &
	微分時間 \({T}_{D}\) \\ \hline
	P 制御 &
	0.5 \({K}_{Pc}\) &
	2.0 \({P}_{Bc}\) &
	--- &
	--- \\ \hline
	PI 制御 &
	0.45 \({K}_{Pc}\) &
	2.2 \({P}_{Bc}\) &
	0.83 \({T}_{c}\) &
	--- \\ \hline
	PID 制御 &
	0.6 \({K}_{Pc}\) &
	1.6 \({P}_{Bc}\) &
	0.5 \({T}_{c}\) &
	0.125 \({T}_{c}\) \\ \hline
\end{tabular}
\end{table}
% /////////////////////////////////////////////////////

\textbf{表~\ref{table:table01}} のように，表番号を参照することもできます．